% ===============================================================
\chapter{Russell's Solutions to Frege's Problems, and Critiques of the Theory of Descriptions}
\chaptermark{Russell's solutions and critiques}
% ===============================================================

\section{The problem of negative existentials}

Assume the sentence ``Pegasus does not exist''. For ``Pegasus'' to have a property, it must be mapped to an object that exists, but the predicate requires that it must not exist, and there seems to be a contradiction.

\section{Collapse of the Fregean Referential Theory of Meaning}

\begin{definition}[Assumption of reference]
All singular terms are assumed to refer.
\end{definition}

When we convert the statement into logical form:
\begin{premises}
  \item[P1)] There does not exist an object such that it is identical to Pegasus.
  \item[P2)] Pegasus is a singular term.
\end{premises}
P1) and P2) are contradictory because of the assumption of reference.

\section{Russell's solution to empty names}

\begin{note}
In all the solutions below, names like ``Pegasus'', ``Venus'' etc.\ are disguised definite descriptions in the view of Russell.
\end{note}

Consider the sentence (a): ``The winged horse of Greek mythology flies''. The following argues that Frege's RTM is incorrect:

\begin{argument}\leavevmode
\begin{premises}
  \item[P1)] RTM is correct \emph{if and only if} sentences with singular terms are meaningful if those sentences have a referent.
  \item[P2)] ``The winged horse of Greek mythology'' is a singular term.
  \item[P3)] ``The winged horse of Greek mythology'' does not have a referent.
  \item[C1)] P1) $\Rightarrow$ (a) is not meaningful.
  \item[P4)] (a) is meaningful.
  \item[C2)] RTM is incorrect, by contradiction.
\end{premises}
\end{argument}

Russell's theory rejects P2) (it is a quantifier instead), and P4) is still true because (a) is assigned a semantic value. Therefore the descriptive theory does hold.

\section{Russell's solution to the problem of substitution}

Why RTM collapses for, for example, a propositional attitude sentence:

\begin{argument}\leavevmode
\begin{premises}
  \item[P1)] Coreferential singular terms ($a$, $b$) have the same semantic value.
  \item[P2)] SSV holds for sentences containing the two terms $a$ and $b$.
  \item[P3)] The definite descriptions of those sentences are coreferential singular terms.
  \item[C1)] So, the sentences should have the same semantic value.
  \item[P4)] They don't.
\end{premises}
\end{argument}

Russell rejects P3), since they are not singular terms.

\section{Russell's solution to the problem of cognitive significance}

Why RTM collapses:

\begin{argument}\leavevmode
\begin{premises}
  \item[P1)] Coreferential singular terms have the same semantic values.
  \item[P2)] Expressions having the same meaning are equally informative.
  \item[C1)] RTM is correct $\Rightarrow$ coreferential singular terms are equally informative.
  \item[P3)] The definite descriptions in the two identity statements $a = a$ and $a = b$ are coreferential singular terms.
  \item[C2)] RTM $\Rightarrow$ $a = a$ and $a = b$ are equally informative.
  \item[P4)] $a = a$ and $a = b$ are not equally informative.
  \item[C3)] RTM is incorrect.
\end{premises}
\end{argument}

Russell rejects P3) because (*drumroll*) they are not singular terms.

\section{Russell's solution to the problem of negative existentials}

Why RTM collapses: let (a) be ``Pegasus exists'' and (b) be ``Pegasus does not exist''.

\begin{argument}\leavevmode
\begin{premises}
  \item[P1)] By the law of excluded middle, only (a) or (b) is true.
  \item[P2)] (a) is true $\Rightarrow$ Pegasus exists.
  \item[P3.0)] The assumption of reference presupposes the following:
  \begin{premises}
    \item[P3)] (b) is true $\Rightarrow$ Pegasus does not exist.
    \item[P3.1)] (b) is true $\Rightarrow$ there exists a Pegasus with the property of not existing.
  \end{premises}
  \item[C1)] Due to the contradiction in P3.1), Pegasus exists.
\end{premises}
\end{argument}

Russell rejects P3.1), not because it is a singular term (surprisingly), but because the predicate ``not existing'' is ambiguous, with the following two meanings when combined with a definite description:
\begin{enumerate}
  \item There is exactly one thing such that it is not the case that it is a winged horse.
  \item It is not the case that there is exactly one thing such that it is a winged horse.
\end{enumerate}
We can pick the second meaning in order to avoid a collapse of the theory of descriptions.

\section{Critiques of Russell}

\subsection*{Communicative success}
\begin{itemize}
  \item Different people might have distinct definite descriptions, and therefore meanings, attributed to the same name. This should imply that communication as a whole is impossible.
  \item The description may not be unique/bijective.
  \item As we will see later, the description might just end up being wrong (``Christopher Columbus'' as the first person to discover America).
\end{itemize}

\subsection*{Distinct logical forms}
Not all logical forms of a sentence are reducible to the same expression. This leads to unsuccessful communication.

\subsection*{Still unable to fully solve the issue of SSV}
\begin{itemize}
  \item Consider our earlier argument, where we said that coreferential singular terms were what led to the collapse of RTM.
  \item Replace that with coextensional predicates (same definite descriptions) and the same collapse will occur for the theory of definite descriptions.
\end{itemize}

\section{Millianism and Russellianism}

\begin{definition}[Millianism]
John Stuart Mill had a philosophy that was often associated with Russellianism. Millianism is essentially just saying that the meaning of a name is its referent.
\end{definition}

Frege disagrees, because meaning according to Frege is sense rather than a referent.

\begin{definition}[Russellianism]
Things that are named by names are constituents of propositions.
\end{definition}

Because these named things are parts of a proposition and are required to deconstruct a proposition, Millianism is often associated with Russellianism. However, Russell is not a Millian, because a Millian believes in singular terms. Russell sees the meaning of names not as singular terms but as definite descriptions.

\section{Scope ambiguities}

\begin{definition}
Sentences which have multiple meanings depending on what kind of scope each function takes are \emph{scope ambiguous}.
\end{definition}

Scope ambiguities involving time and negation are in accordance with the tense of the sentence in a Russellian analysis. For example, a ``present'' king of France vs a ``past'' king of France being bald changes the semantic value of the sentence entirely.
